% Bagian: intuitionistic-logic
% Bab: semantics
% Subbagian: introduction

\documentclass[../../../include/open-logic-section]{subfiles}

\begin{document}

\olfileid[id]{int}{sem}{int}

\olsection{Pendahuluan}

Tidak ada logika yang diuraikan secara memuaskan tanpa semantik, dan logika
intuisionistik bukanlah pengecualian. Dalam logika klasik, semantik berdasarkan
!!{valuation}s bersifat kanonis, sedangkan logika intuisionistik mempunyai
beberapa semantik yang saling bersaing. Tidak satu pun sepenuhnya memuaskan
dalam arti memberikan penjelasan tentang makna konektif yang dapat diterima
secara intuisionistik.

Semantik berdasarkan !!{relational model}s, yang serupa dengan semantik bagi
logika modal, mungkin merupakan semantik yang paling populer. Dalam semantik
ini, !!{propositional variable}s ditetapkan kepada dunia-dunia, dan dunia-dunia
tersebut dihubungkan oleh relasi keterjangkauan. Relasi itu selalu merupakan
urutan parsial; yakni, relasi itu refleksif, antisimetris, dan transitif.

Secara intuitif, dunia-dunia ini dapat dipandang sebagai keadaan pengetahuan
atau ``situasi evidensial''. Suatu keadaan~$w'$ terjangkau dari~$w$ jika dan
hanya jika, sejauh yang kita ketahui, $w'$ merupakan keadaan pengetahuan
(masa depan) yang mungkin; yakni, keadaan yang selaras dengan apa yang
diketahui pada~$w$. Setelah suatu proposisi diketahui, pengetahuan itu tidak
dapat hilang; yakni, jika $!A$ diketahui pada~$w$ dan~$Rww'$, maka $!A$ juga
diketahui pada~$w'$. Jadi, ``pengetahuan'' bersifat monoton terhadap relasi
keterjangkauan.

Jika kita mendefinisikan ``$!A$ diketahui'', sebagaimana dalam logika
epistemik, sebagai ``benar di semua alternatif epistemik'', maka $!A \land !B$
diketahui pada~$w$ jika di semua alternatif epistemik, baik $!A$ maupun~$!B$
diketahui. Namun, karena pengetahuan bersifat monoton dan $R$ refleksif, hal
ini berarti bahwa $!A \land !B$ diketahui pada $w$ jika dan hanya jika $!A$
dan $!B$ diketahui pada~$w$. Dengan alasan yang sama, $!A \lor !B$ diketahui
pada $w$ jika dan hanya jika sekurang-kurangnya salah satunya diketahui. Jadi,
bagi $\land$ dan $\lor$, syarat kebenaran konektif tersebut bertepatan dengan
syarat dalam logika klasik.

Namun, syarat kebenaran bagi kondisional berbeda dari syarat dalam logika
klasik. $!A \lif !B$ diketahui pada~$w$ jika dan hanya jika tidak ada $w'$
dengan $Rww'$ tempat $!A$ diketahui tanpa~$!B$ juga diketahui. Syarat ini tidak
sama dengan syarat bahwa $!A$ tidak diketahui atau $!B$ diketahui pada~$w$.
Sebab, jika baik $!A$ maupun $!B$ tidak kita ketahui pada~$w$, mungkin ada
keadaan epistemik masa depan~$w'$ dengan $Rww'$ sedemikian sehingga pada $w'$,
$!A$ diketahui tanpa~$!B$ juga diketahui.

Kita mengetahui $\lnot !A$ hanya jika tidak ada keadaan epistemik masa depan
yang mungkin tempat kita mengetahui~$!A$. Gagasannya ialah bahwa jika $!A$
dapat diketahui, maka pada suatu keadaan epistemik masa depan yang mungkin,
$!A$ menjadi diketahui. Karena kita tidak dapat mengetahui $\lfalse$, pada
keadaan epistemik masa depan itu kita akan mengetahui $!A$, tetapi tidak
mengetahui~$\lfalse$.

Dalam interpretasi ini, prinsip tiada jalan tengah gagal. Sebab, ada beberapa
$!A$ yang belum kita ketahui, tetapi mungkin kelak kita ketahui. Bagi
!!a{formula}~$!A$ semacam itu, baik $!A$ maupun~$\lnot !A$ tidak diketahui,
sehingga $!A \lor \lnot !A$ tidak diketahui. Namun, misalnya, kita mengetahui
bahwa $\lnot(!A \land \lnot !A)$. Tidak ada keadaan masa depan yang mungkin
tempat kita mengetahui baik $!A$ maupun $\lnot !A$, dan hal ini kita ketahui
secara terlepas dari apakah kita mengetahui~$!A$ atau $\lnot !A$.

!!^{relational model}s bukanlah satu-satunya semantik yang tersedia bagi
logika intuisionistik. Semantik topologis merupakan pilihan lain: dalam
semantik ini, proposisi ditafsirkan sebagai himpunan terbuka dalam suatu ruang
topologis, dan konektif ditafsirkan sebagai operasi pada himpunan-himpunan
tersebut (misalnya, $\land$ bersesuaian dengan irisan).

\end{document}
